\section{Metodologia}

A Figura~\ref{fig:pipeline} sintetiza o desenho experimental. A partir de uma única versão auditada do TACO, a partição e as caixas foram mantidas constantes, três taxonomias foram derivadas e treinadas sob protocolo idêntico, e cinco análises complementares foram aplicadas ao mesmo conjunto de teste.

\begin{figure}[!t]
  \centering
  \includegraphics[width=.97\linewidth]{methodology_pipeline.pdf}
  \caption{Fluxo experimental. Somente a taxonomia varia entre os três ramos de treinamento; imagens, caixas, partição, arquitetura e hiperparâmetros permanecem fixos.}
  \label{fig:pipeline}
\end{figure}

\subsection{Dados, auditoria e partição}

Foi utilizado exclusivamente o arquivo \path{data/annotations.json}, obtido do repositório oficial TACO \cite{proenca2020tacorepo}. As anotações seguem o formato COCO \cite{lin2014coco}; somente as caixas delimitadoras (\textit{bounding boxes}) foram consideradas. A auditoria programática confirmou 1.500 imagens, 4.784 objetos, 60 categorias, nenhuma imagem vazia e nenhum arquivo ausente. Sete caixas excediam a borda em no máximo 1,32 pixel; o recorte foi efetuado apenas durante a conversão, preservando-se o JSON oficial e registrando-se cada alteração.

A partição foi sorteada uma vez por imagem com semente 42: treino, validação e teste receberam, respectivamente, 1.050/225/225 imagens e 3.359/788/637 objetos. As três taxonomias reutilizam exatamente essa partição. As coordenadas COCO $(x,y,w,h)$ foram convertidas para centro, largura e altura normalizados do YOLO; scripts verificaram limites, dimensões positivas, IDs e igualdade das partições. Imagens com orientação EXIF foram transpostas fisicamente antes da conversão.

\begin{figure}[H]
  \centering
  \includegraphics[width=.94\linewidth]{class_rank_distribution.pdf}
  \caption{Instâncias de treino das 60 classes Fine, ordenadas por frequência. As linhas marcam os limiares da análise de cauda longa.}
  \label{fig:distribution}
\end{figure}

\subsection{Taxonomias}

Fine preserva as 60 categorias oficiais. Material remapeia todas elas para Plastic (24 categorias de origem), Metal (9), Glass (4), Paper/Cardboard (13), Organic (1) e Other (9). Compostos ou itens sem associação natural, como bateria, sapato e embalagens multicamadas, foram atribuídos a Other; cada decisão e sua justificativa estão em \path{configs/material_mapping.json}. Binary atribui \textit{Litter} a todas as instâncias. As caixas permanecem inalteradas entre as versões. A Figura~\ref{fig:distribution} evidencia a cauda longa que motiva os agrupamentos.

\subsection{Modelo e protocolo experimental}

O YOLO26n foi inicializado com pesos COCO e treinado com entrada $640\times640$, lote 24, até 60 épocas, parada antecipada (\textit{early stopping}) com paciência 15, AdamW \cite{loshchilov2019adamw} e taxa inicial $10^{-3}$. Foram mantidas as transformações padrão da implementação Ultralytics, a precisão mista e a execução determinística. Os experimentos utilizaram Python 3.12.10, PyTorch 2.8.0+cu128 e Ultralytics 8.4.157 em uma RTX 3060 Laptop de 6~GB. As sementes 42, 123 e 2026 foram aplicadas às três taxonomias, sem ajuste específico por configuração.

\subsection{Avaliação e incerteza}

Todas as análises quantitativas reportadas empregam o mesmo avaliador controlado do projeto: comparação dos modelos treinados, reavaliação hierárquica e bootstrap. As predições dos nove modelos foram reavaliadas sem novo treinamento ou inferência, verificando-se igualdade numérica entre o resultado Fine da tabela principal e sua referência hierárquica em cada semente, com tolerância $10^{-12}$. As saídas foram exportadas com confiança mínima 0,001 e limite de 300 detecções por imagem. Os resultados nativos Ultralytics permanecem arquivados, mas não integram as tabelas ou figuras de desempenho do artigo. A tabela principal apresenta média e desvio-padrão entre sementes.

Para uma caixa predita $B$ e uma verdade-terreno $G$, a sobreposição é definida por
\begin{equation}
  \operatorname{IoU}(B,G)=\frac{|B\cap G|}{|B\cup G|}.
  \label{eq:iou}
\end{equation}
Em cada limiar $\tau$, o AP de cada classe é calculado como a média do envelope de precisão em 101 níveis de revocação; a métrica reportada é a macro-média das classes com instâncias no teste. São exigidas classe correta e associação espacial um-para-um. A métrica mais estrita agrega dez limiares de IoU, de 0,50 a 0,95:
\begin{equation}
  \operatorname{AP}_{50:95}=\frac{1}{10}\sum_{k=0}^{9}
  \operatorname{AP}_{0{,}50+0{,}05k}.
  \label{eq:map}
\end{equation}
Para cada classe, o avaliador seleciona o ponto de confiança que maximiza $F_1=2PR/(P+R)$ em IoU de 0,50. Precisão, revocação e $F_1$ são então agregados por macro-média. Assim, o $F_1$ reportado não corresponde a um único limiar global nem à média harmônica das macro-médias de precisão e revocação.

Para decompor o erro, um avaliador controlado faz casamento guloso um-para-um em ordem decrescente de confiança. Seja $\pi_L:\mathcal{C}_{Fine}\rightarrow\mathcal{C}_L$ o mapeamento para o nível $L\in\{Fine,Material,Binary\}$. Cada detecção $(B_i,c_i,s_i)$ é reavaliada como
\begin{equation}
  (B_i,c_i,s_i)\mapsto(B_i,\pi_L(c_i),s_i),
  \label{eq:remap}
\end{equation}
preservando geometria e confiança. A linha Fine é a mesma referência da comparação dos modelos treinados. Como o número de classes e a macro-média também mudam, o ganho hierárquico combina o crédito recuperado por confusões colapsadas com o efeito da agregação de categorias.

Foi aplicado bootstrap percentil pareado \cite{efron1993bootstrap} para a obtenção de ICs de 95\%. Em cada réplica $b$, as 225 imagens são reamostradas com reposição e calcula-se
\begin{equation}
  \Delta^{(b)}=M_A(\mathcal{I}^{(b)})-M_B(\mathcal{I}^{(b)}),\quad
  IC_{95\%}=[Q_{0{,}025}(\Delta),Q_{0{,}975}(\Delta)].
  \label{eq:bootstrap}
\end{equation}
Foram realizadas 2.000 reamostragens para cada uma das três sementes, com índices idênticos nos dois termos de cada comparação. São apresentados a média e o desvio-padrão das diferenças médias do bootstrap entre sementes. Para uniformizar a interpretação, as diferenças foram orientadas de modo que valores positivos favoreçam a taxonomia mais simples; também se contabilizou o número de sementes cujo IC excluiu zero. O cálculo foi acelerado mediante pré-computação dos casamentos por classe e limiar de IoU; testes automatizados confirmaram equivalência numérica com o avaliador original a $10^{-12}$. Para a análise de cauda longa, $n$ corresponde ao número de instâncias de treino por classe: rara ($n\leq19$), média ($20\leq n\leq100$) e frequente ($n>100$). Esses cortes constituem estratos operacionais deste estudo, e não uma definição universal de cauda longa. As métricas por classe incluem somente rótulos presentes no teste fixo.
